\chapter{模型简介}
\label{chap:machine-learning-models}

有了带标签的历史邮件，我们仍有几种不同的办法判断新邮件：学习一组特征权重，建立一棵
判断规则树，或者保存历史邮件，在预测时寻找相似样本。第\ref{chap:introduction}章的邮件
案例展示了根据标签调整权重的过程，但同样的标签也能用于学习其他形式的规律。怎样表示
这些规律，是选择模型时需要回答的问题。

怎样比较这些模型？我们需要知道，它们要保存多少信息，能否说明预测的不确定性，又根据
什么关系把输入变成结果。这些差别还会影响训练时如何评价预测。本章从邮件例子出发，
区分候选模型、具体模型与训练过程，再把这些问题整理成五个观察维度。分类与聚类的对照
将帮助我们看清：两个方法在哪里采用了相同假设，在哪里作出了不同选择。

\section{什么是模型}
\label{sec:models-meaning}

继续使用第\ref{chap:introduction}章的邮件例子。每封邮件由促销提示词次数与外部链接数
表示，记为$\bm{x}=(x_1,x_2)^{\top}$。规定“对这两个特征加权求和，再转成垃圾邮件的
概率”，只是选定了一种预测形式。只要权重还没有确定，同一封邮件就可能得到不同的概率。
这里需要区分两个对象：训练前允许选择的整组候选关系，以及训练后实际使用的某一个关系。
满足给定形式与约束的一组候选模型称为\term{模型族}（model family），记为$\mathcal H$；
其中一个具体模型记为$h\in\mathcal H$。预测任务中的模型族通常也称为假设空间；“模型族”
则还适用于聚类、降维和分布建模等不只输出预测标签的任务。

模型族给出选择范围，具体模型给出计算依据。邮件分类需要确定的量是各特征的权重；换成
其他任务，也可能要确定代表群组的中心位置，或压缩数据时使用的投影方向。

在给定表示方式下，需要从数据中估计、用于确定具体模型的量是模型参数。回归系数、聚类
中心的坐标和神经网络权重都属于这类参数。用$\bm{\theta}$统称这些量时，
$h_{\bm{\theta}}$表示由这组参数确定的模型。

沿用前文的邮件预测式，用$\bm{\theta}=(b,w_1,w_2)^{\top}$记录参数，可以把具体模型
与模型族写在一起：
\begin{equation}
  \begin{aligned}
    h_{\bm{\theta}}(\bm{x})
      &=\frac{1}{1+\exp[-(b+w_1x_1+w_2x_2)]},\\
    \mathcal{H}&=\{h_{\bm{\theta}}:\bm{\theta}\in\mathbb{R}^{3}\}.
  \end{aligned}
  \label{eq:models-family-example}
\end{equation}
其中，$w_1,w_2$是两个特征的权重，$b$是在加权和之外加入的偏置。
集合$\mathcal{H}$包含所有允许的权重组合所确定的模型；例如，令
$b=-2$、$w_1=0.8$、$w_2=0.3$，就选出了其中一个具体模型。
这里的数值仅用于说明概念。换一组权重，通常会得到同一模型族中的另一个模型；
把加权和换成条件分支树，则改变了所考虑的模型族。

确定候选范围以后，还要有一种办法根据数据作出选择，这才是学习算法的工作。在邮件
例子中，算法根据训练样本调整三个参数。换一种优化算法，调整参数的步骤可能不同，
允许选择的预测关系却可以保持不变。模型族、模型参数和学习算法分别说明“可以选什么”
“选中了什么”和“怎样选”。

除了要估计的模型参数，还常需指定超参数，例如网络的层数、聚类的簇数、近邻方法的邻居数。
超参数规定候选范围或学习规则，也可以借助验证数据选择。例如，先确定网络有几层，
才能确定需要学习哪些连接权重；先确定每次参考几个近邻，才能执行相应的投票规则。
一个方法的超参数很少，预测时仍可能需要保存大量信息，下一节的近邻方法就属于这种情形。

邮件也可以提出分类之外的问题。若想把相似邮件放在一起，可以用若干代表点描述群组；
若想压缩特征，可以学习一个低维子空间；若想描述邮件特征怎样共同出现，可以建立它们的
概率分布。模型因任务而改变，输出可能是一个标签、一组坐标或一个分布，但都需要说明
保存了什么规律、怎样使用这些规律，以及怎样从数据中得到它们。

表\ref{tab:models-dimensions}列出本章采用的五个观察维度。这是一种组织和比较模型的方法，
并非把所有模型放进一棵互不重叠的分类树。参数化方式与表示规模、概率语义和关系形式
分别考察表示规模、分布假设与变量关系；输出接口考察模型在使用时返回什么；
损失与学习准则则考察学习时用什么准则评价候选模型。

\begin{table}[htbp]
  \centering
  \small
  \begin{tabularx}{\linewidth}{@{}lXX@{}}
    \toprule
    \textbf{维度} & \textbf{要回答的问题} & \textbf{主要观察方向} \\
    \midrule
参数化方式与表示规模 & 模型表示规模是否有固定上限？ & 参数化、非参数化及半参数化组合 \\
    概率语义 & 是否为有关变量明确建立概率分布？ & 非概率、概率 \\
    关系形式 & 有关变量怎样发生联系？ & 线性、非线性 \\
    输出接口 & 使用模型时返回什么？ & 组织形式与结果含义 \\
    损失与学习准则 & 怎样评价候选模型？ & 逐样本差异、关系与结构、轨迹与长期目标 \\
    \bottomrule
  \end{tabularx}
  \caption{认识模型的五个维度。每个维度回答不同问题，不能由其中一项直接推断其余各项。}
  \label{tab:models-dimensions}
\end{table}

式\eqref{eq:models-family-example}给出的逻辑回归就同时涉及这五个问题：它保存三个参数，
通过加权和计算类别概率，并根据概率与已知标签的吻合程度学习参数。如果只看最终返回的
“垃圾邮件”标签，这些差别就被隐藏了。下面从模型需要保存多少信息入手，逐步展开这些
问题，并用预测、聚类与压缩中的例子作比较。

\section{参数化方式与表示规模}
\label{sec:models-capacity}

训练邮件从一百封增加到一万封后，模型要保存的信息会怎样变化？有的方法仍只更新原来的
几个权重，有的方法则保留更多样本或增加更多判断分支。这种差别是理解\term{容量}
（capacity）的一个起点：容量关心一个模型族能够表达多丰富的关系。

本节考察表示规模是否受预先固定的参数维数限制，并介绍两类表示的组合。这里观察的是
整个候选范围：一次训练得到的模型即使能用有限信息保存，也不能据此断定所有候选模型
的表示规模都有固定上限。参数范围和结构限制同样会影响容量；它们与泛化的理论联系
留到第\ref{part:learning-theory}部分。

\subsection{参数化模型}
\label{sec:models-parametric}

固定输入表示与模型配置后，如果模型族中的候选模型由预先固定维数的参数向量描述，
就称为\term{参数化模型族}（parametric model family），通常简称参数化模型。
例如，将允许的参数向量限制在$\mathbb{R}^{p}$中的某个范围内，$p$便是参数维数。
这里的关键是$p$不会仅因训练样本数$n$增加而自动增长，而参数的具体取值仍要从数据中估计。

以$d$维输入的线性回归为例，每个模型由$d$个权重与一个偏置确定，所以$p=d+1$。
当训练数据从一百个样本增加到一万个时，需要重新估计的数值可能改变，
但仍是这$d+1$个数值。训练以后，只用这些权重与偏置就能对新输入作预测。
固定层数、宽度与连接方式的神经网络也有这样的性质；其参数可以很多，
但数量由配置确定，不随样本数量自动增加。

参数化模型并不限于预测任务。用于聚类的\term{K均值模型}（K-means）用$k$个中心代表
数据中的群组，把新样本分给最近的中心。这类代表点也称为原型。
固定输入维数$d$和簇数$k$后，模型参数是$k$个中心的坐标，共$kd$个数值。
例如，二维数据分成两簇时，只需保存两个二维中心，即四个数值，再配合最近中心规则
完成分配。增加训练样本可以改变中心位置，但无需为每个新输入增添一组全局模型参数。

在表示学习中，\term{主成分分析}（principal component analysis, PCA）用少量方向
表示数据的主要变化。
固定原始维数与保留方向数以后，均值和投影方向也可以用固定规模的数值表示。
这里的簇数、保留方向数以及网络宽度是模型的超参数。
在确定一个模型是参数化还是非参数化时，要先固定这些配置，再观察模型族的表示规模如何随数据变化。

\subsection{非参数化模型}
\label{sec:models-nonparametric}

如果希望在数据增加时保留更细的局部关系，就可以允许表示规模一同增长。这类方法所用的
\term{非参数化模型族}（nonparametric model family），不把候选关系限制在预先固定的
有限维参数形式中。“非参数化”描述的是这个限制被放宽了，模型仍然有待确定的数值，
算法也仍有需要选择的设置。

\term{$k$近邻}（$k$-nearest neighbors, $k$NN）分类器采用了一种直接的做法：保存训练
样本及其标签，预测新邮件时找出特征最接近的$k$封历史邮件，再按标签投票。
这里的$k$控制每次查询所用的邻居数量，属于上一节所说的超参数。

只知道$k=3$，还无法对一封新邮件分类：必须知道保存了哪些样本、它们位于什么位置，
以及它们的标签。同样设置$k=3$，两份不同的训练数据可以给出不同的预测规则。
标准kNN在一百个样本上保存一百组特征和标签，在一万个样本上则保存一万组；
每次查询仍只取三个邻居，但具体是哪三个，要由新输入与全部已保存样本之间的关系决定。

预测仍依赖全部已保存实例，而不只是数字$3$。随着训练样本增加，不同位置附近可用于
预测的信息增多，模型能够形成更细的局部划分。这种表示方式不受固定维数的系数向量限制，
是kNN被归为非参数化方法的原因。即使用验证集选择$k$，$k$控制邻域范围、实例提供
预测依据的分工也没有改变。

这与K-means中的$k$作用不同。固定K-means的簇数，就固定了需要学习的中心数量；
固定kNN的邻居数，只固定一次查询的参与者数量，并未固定需要保留的实例数量。
图\ref{fig:models-neighbors-prototypes}用这两种方法对照超参数与模型表示之间的关系。

\begin{figure}[htbp]
  \centering
  % 原创教学示意；近邻数与簇数是不同超参数。
% 左图查询点 (15,13) 的三个最近邻为 (13,15)、(17,11)、(20,16)。
\begin{tikzpicture}[x=1mm,y=1mm,font=\footnotesize]
  \foreach \shift in {0,54}{
    \begin{scope}[xshift=\shift mm]
      \fill[BookPaper] (0,0) rectangle (46,35);
      \foreach \x/\y in {8/8,13/15,17/11,10/19,20/16}{
        \fill[FigureInput] (\x,\y) circle (0.9);
      }
      \foreach \x/\y in {29/23,34/29,38/23,30/30,39/31}{
        \fill[FigureLoss] (\x-0.8,\y-0.8) rectangle (\x+0.8,\y+0.8);
      }
    \end{scope}
  }
  \node[font=\heiti\small] at (23,41) {$k$近邻：$k=3$};
  \node[font=\heiti\small] at (77,41) {K-means：$k=2$};
  \draw[FigureModel,dashed,thick] (15,13) circle (6.7);
  \fill[BookInk] (15,14.2)--(16.2,13)--(15,11.8)--(13.8,13)--cycle;
  \node[align=center] at (23,-7) {保存样本及标签\\每次仅选近邻投票};
  \begin{scope}[xshift=54mm]
    \fill[BookPaper,opacity=0.82] (0,0) rectangle (46,35);
    \foreach \x/\y in {13.6/13.8,34/27.2}{
      \draw[FigureModel,very thick] (\x-2,\y)--(\x+2,\y) (\x,\y-2)--(\x,\y+2);
    }
  \end{scope}
  \node[align=center] at (77,-7) {保存两个中心\\按最近中心分配};
\end{tikzpicture}

  \caption{两个$k$控制不同的对象。左侧kNN保存十个带标签样本，菱形表示新输入，
  虚线圈内的三个近邻参加投票；右侧K-means用两个中心表示群组。
  固定近邻数不会固定实例总数，固定簇数则固定了中心数量。图为表示方式的教学示意。}
  \label{fig:models-neighbors-prototypes}
\end{figure}

非参数化也不意味着无需假设。kNN依赖所选距离、特征尺度和邻居数：
距离决定哪些样本接近，$k$决定一次投票参考多大的局部范围。
保存更多实例会增加存储与查询成本，改变特征尺度也可能改变邻居关系。
这里保存实例是为了完成预测。线性回归的程序也可能为重新训练、检查误差或记录来源而
保存全部数据，但其预测关系仍能由固定数量的权重表示。区分两类模型，需要看预测关系
怎样表示；程序中恰好保留了多少文件或样本，本身不足以作判断。

局部关系也可以通过增加规则来细化。\term{决策树}（decision tree）根据特征阈值逐步
分支，把输入分到不同区域，再给出各区域的预测。
如果不预先固定叶节点数的上限，更多样本可以支持更多分支，所需阈值和叶节点输出的
数量也可以增加，因此属于非参数化设定。若将树深度或叶节点数固定在某个上限内，
则变成了表示规模有固定上限的设定。树的区域划分将在第\ref{sec:models-nonlinear}节说明。

\term{随机森林}（random forest）综合多棵决策树的预测。
即使树的棵数固定，只要各树的叶节点数量仍可随数据增长，整体也可以采用非参数化设定。
若棵数与每棵树的叶节点上限都固定，表示规模才相应受到固定限制。
这与kNN类似：固定一个超参数，不一定固定模型的全部自由度。

\term{核回归}（kernel regression）按输入之间的距离或相似性，对训练样本的目标值
作局部加权平均。控制邻近范围的尺度通常称为带宽，它是超参数；
预测所依据的样本位置和目标值仍可随数据增加。
例如，带宽不变时，增加附近的房屋成交记录仍会改变当前位置的加权房价预测。
这种由实例不断补充局部关系的方式，也属于非参数化方法。
近邻、树与局部平滑方法的联系可在统计学习教材中进一步比较\scite{hastie2009}

\subsection{半参数化模型}
\label{sec:models-semiparametric}

用面积与楼龄预测房价时，我们可能愿意假设面积的作用是线性的，却不确定楼龄与房价之间
适合用什么曲线。此时可以保留明确的面积项，让楼龄项从数据中更灵活地学习。
\term{半参数化模型}（semiparametric model）就把固定有限维的参数化部分与非参数化
部分结合起来。记面积为$x$、楼龄为$z$，一个教学设定是
\begin{equation}
  f_{\beta,g}(x,z)=\beta x+g(z).
  \label{eq:models-semiparametric}
\end{equation}
系数$\beta$描述面积的固定线性作用，未知函数$g$描述楼龄与房价之间的关系，
后者可以用局部平滑等非参数化方法估计。学习时既要确定$\beta$，也要估计整条函数$g$。
“半参数化”指这两部分的结合，并不是参数数量介于某两个模型之间，也不是一半参数已知。

同一种方法也可能采用不同的容量设定。以\term{支持向量机}（support vector machine, SVM）
为例，它根据得分判断类别，学习时要求样本尽量位于边界的正确一侧，并与边界保持一定
间隔。这里考察其表示规模，间隔目标将在
第\ref{sec:models-spatial-objectives}节介绍。

固定$d$维输入的线性SVM使用$s(\bm{x})=\bm{w}^{\top}\bm{x}+b$，
全部预测关系由$d$个权重与一个偏置确定，所以属于参数化设定。
采用核函数时，得分则常写成训练样本贡献的加权和：
\begin{equation}
  s(\bm{x})=\sum_{i\in I_{\mathrm{SV}}}
    a_i\,\kappa(\bm{x}_i,\bm{x})+b.
  \label{eq:models-support-vector-expansion}
\end{equation}
其中$\kappa$是用于计算两个输入之间关系的核函数，$a_i$为学习得到的系数，
$I_{\mathrm{SV}}\subseteq\{1,\ldots,n\}$包含非零系数对应的训练样本索引，
这些样本称为\term{支持向量}（support vectors）。核函数与特征变换的关系见
第\ref{sec:models-nonlinear}节。

对于通常的高斯核SVM，支持向量的数量可以随样本量增加，模型也没有等价的固定有限维
特征表示，因此通常归为非参数化模型。只保留一部分训练样本，可以使具体模型更紧凑，
但不意味着整个模型族变成半参数化；固定核的带宽或其他训练超参数，也不会固定支持向量数。
如果核对应固定的有限维特征，例如固定次数的多项式核，则又可以用固定维数的权重表示，
需要按这一具体设定判断。

要得到半参数化SVM，需要在预测关系中明确结合两种部分。例如，让得分或回归预测包含
一个可解释的线性趋势，再加上由非参数化核函数模型学习的修正项；前者具有固定数量的
系数，后者可以随数据扩展。
已有研究将支持向量方法扩展到这种组合形式，并联合学习两部分\scite{smola1998semiparametric}
SVM的容量取决于这些具体设定：线性形式具有固定维数的表示；核形式要看核对应的特征
空间；半参数化形式则明确组合了两类部分。超参数只有几个，或某次训练得到的支持向量
只有有限个，都不能替代对整个候选范围的判断。

\section{变量的概率语义}
\label{sec:models-probability}

即使确定了模型的表示规模，预测时仍有一个选择：只给出结果，还是同时说明结果可能怎样
变化？假设根据今天的气象观测预测明天的温度，模型可以直接给出$22\,^{\circ}\mathrm{C}$，
也可以描述温度落在不同范围内的可能性，例如主要集中在$20$到$24\,^{\circ}\mathrm{C}$之间。
后一种做法需要为可能的结果建立概率分布。这个选择同样会出现在聚类、降维与分类中。

\subsection{非概率模型}
\label{sec:models-nonprobabilistic}

\term{非概率模型}（non-probabilistic model）不为所考察的变量明确建立概率分布，
而是用函数、规则或几何关系直接确定结果。
在温度预测中，记气象特征为$\bm{x}$、模型参数为$\bm{\theta}$，模型计算
\begin{equation}
  \hat{y}=f_{\bm{\theta}}(\bm{x}).
  \label{eq:models-point-prediction}
\end{equation}
训练时，根据预测温度与实际温度的偏差调整参数；使用时，将新的观测代入函数，得到一个
预测值。例如，输出$22\,^{\circ}\mathrm{C}$表示模型给出的温度预测。这个数值本身没有描述
预测可能偏离多少，也没有说明较低或较高温度出现的可能程度。

K-means和PCA也可以直接用几何关系定义。K-means寻找若干中心，使样本尽量靠近各自的
中心；给定新样本后，比较距离并返回最近中心的编号。PCA寻找少量投影方向，用低维坐标
保留数据的主要变化；给定新样本后，计算它沿这些方向的坐标。两者分别依赖距离与投影，
不需要先规定样本服从什么概率分布。

温度误差、簇内距离和重构误差都能直接用来评价这些模型。当任务需要的是一个预测值、
一个分组或一组坐标时，这样的建模已经可以完成工作。如果还要判断温度偏离预测值超过
两度的概率有多大，就需要额外的概率描述。

\subsection{概率模型}
\label{sec:models-probabilistic}

同样的气象特征未必对应完全相同的明日温度。要保留这种变化，可以把温度看成一个随机
变量，用分布描述各个范围内出现的可能性。\term{概率模型}（probabilistic model）就是
为任务涉及的随机变量及其关系建立概率分布的模型。对于温度预测，可以让函数给出分布的
均值，并用方差描述围绕均值的变化：
\begin{equation}
  y\mid\bm{x}\sim\mathcal{N}\bigl(\mu_{\bm{\theta}}(\bm{x}),\sigma^2\bigr).
  \label{eq:models-linear-gaussian}
\end{equation}
这里$y\mid\bm{x}$表示给定气象特征$\bm{x}$后的温度，$\mu_{\bm{\theta}}(\bm{x})$给出
该条件下的均值；$\mathcal{N}(\mu,\sigma^2)$表示均值为$\mu$、方差为$\sigma^2>0$的高斯分布。
若均值为$22\,^{\circ}\mathrm{C}$、标准差为$2\,^{\circ}\mathrm{C}$，模型既保留了一个中心预测，
也描述了其他温度的可能程度。图\ref{fig:models-point-distribution}对照了这两种预测形式。

\begin{figure}[htbp]
  \centering
  % 原创教学示意：同一温度预测；右图精确绘制 N(22, 2^2) 的密度。
\begin{tikzpicture}[x=1mm,y=1mm,font=\footnotesize]
  \node[font=\heiti\small] at (22,35) {一个预测值};
  \node[font=\heiti\small] at (80,35) {一个预测分布};
  \begin{scope}
    \draw[book arrow] (0,0)--(45,0);
    \foreach \t/\x in {16/0,20/13.333,22/20,24/26.667,28/40}{
      \draw[BookMuted] (\x,0)--(\x,-1);
      \node[below] at (\x,-1) {$\t$};
    }
    \fill[FigureModel] (20,0) circle (1.5);
    \node[above,text=FigureModel] at (20,6) {$\hat y=22$};
    \draw[FigureModel,dashed] (20,5)--(20,2);
    \node at (22,-11) {温度（$^{\circ}\mathrm{C}$）};
  \end{scope}
  \begin{scope}[xshift=60mm]
    \fill[FigureModel!18] (13.333,0)
      -- plot[domain=13.333:26.667,samples=60] (\x,{19.947*exp(-((\x-20)/6.667)^2/2)})
      --(26.667,0)--cycle;
    \draw[FigureModel,thick]
      plot[domain=0:40,samples=100] (\x,{19.947*exp(-((\x-20)/6.667)^2/2)});
    \draw[book arrow] (0,0)--(44,0);
    \draw[book arrow] (0,0)--(0,26);
    \node[above right,inner sep=1pt] at (0,26) {密度（$(^{\circ}\mathrm{C})^{-1}$）};
    \foreach \y/\lab in {10/0.1,20/0.2}{
      \draw[BookMuted] (-1,\y)--(0,\y);
      \node[left,inner sep=2pt] at (-1,\y) {$\lab$};
    }
    \foreach \t/\x in {16/0,20/13.333,24/26.667,28/40}{
      \draw[BookMuted] (\x,0)--(\x,-1);
      \node[below] at (\x,-1) {$\t$};
    }
    \node at (20,-11) {温度（$^{\circ}\mathrm{C}$）};
  \end{scope}
\end{tikzpicture}

  \caption{同一温度预测的两种描述。左侧只给出一个预测值；右侧给出高斯预测分布，
  阴影面积表示温度落在$20$到$24\,^{\circ}\mathrm{C}$之间的概率。数值用于教学示意。}
  \label{fig:models-point-distribution}
\end{figure}

概率建模并不限于监督预测，聚类也有相应的两种做法。K-means用中心与距离描述群组；
\term{高斯混合模型}
（Gaussian mixture model, GMM）则用若干高斯分布描述数据，每个群组除了中心，还有
各方向上的分散程度以及在总体中所占的比例。对处在两组之间的样本，K-means按最近中心
给出一个编号，GMM可以计算它分别属于两组的概率。这样，群组的大小、形状和重叠程度
都能参与分配判断。

降维也可以作同样的对照。PCA与\term{概率主成分分析}（probabilistic PCA, PPCA）
都寻找低维表示，但二者对数据作出的说明不同。
PCA用线性投影压缩和重构数据；PPCA在低维线性表示的基础上增加高斯分布与观测噪声假设，
从而描述数据在这一低维空间附近怎样分布。它不仅能讨论投影位置，还能在模型假设下估计缺失
观测的分布。两者联系紧密，但PPCA额外回答了“哪些观测更可能出现”这个问题\scite{tipping1999}

\paragraph{分类问题的概率建模选择}

回到邮件分类，使用者关心的是：看到这些特征以后，这封邮件属于各类的概率是多少？
这由类别的条件分布$p(y\mid\bm{x})$回答。得到它有两条路线：直接学习特征与类别概率
之间的关系，或先描述各类邮件的特征分布，再反推类别。两者分别称为判别式建模与
生成式建模。

\term{判别式建模}（discriminative modeling）直接学习输入到类别的关系。
在这里讨论的概率分类中，它直接建立$p(y\mid\bm{x})$。例如，逻辑回归根据带标签样本
调整权重，使给定邮件特征后的类别概率与标签相符。模型不必同时说明各种特征组合在全部邮件中
出现的概率，也不必分别描述正常邮件与垃圾邮件内部的特征分布。
更广义的判别式建模也可以只学习得分或决策规则，例如SVM；本小节关注的是其中直接学习
类别概率的情形。

\term{生成式建模}（generative modeling）用于分类时，先建立输入与类别的联合分布
$p(\bm{x},y)=p(y)p(\bm{x}\mid y)$。其中$p(y)$描述各类别的比例，
$p(\bm{x}\mid y)$描述每一类内部的特征分布。训练得到这些分布后，再根据新输入的特征
反推其类别。对于编号为$1,\ldots,k$的类别，由贝叶斯公式得到
\begin{equation}
  p(y=c\mid\bm{x})
  =\frac{p(y=c)\,p(\bm{x}\mid y=c)}
         {\sum_{j=1}^{k}p(y=j)\,p(\bm{x}\mid y=j)}.
  \label{eq:models-generative-classification}
\end{equation}
其中$c$是待考察类别，分母对所有类别求和，并要求其大于零。分子同时考虑类别原本有多
常见，以及当前特征在该类别下有多可能出现；除以所有类别的总和，就得到观察特征后的
类别概率，也称为类别后验概率。若只需输出概率最大的类别，
比较各类分子即可，因为分母相同。

作为一个数值例子，约定邮件标签$1$为垃圾邮件、$0$为正常邮件。假设模型得到
$p(y=1)=0.2$、$p(y=0)=0.8$；对于某一组计数特征$\bm{x}$，
$p(\bm{x}\mid y=1)=0.12$、$p(\bm{x}\mid y=0)=0.01$。
这些数值仅用于演示计算。两类的联合概率分别为$0.024$和$0.008$，
因此垃圾邮件的后验概率为$0.024/(0.024+0.008)=0.75$。
这个结果既利用了特征的区分能力，也计入了两类邮件的总体比例。

\term{朴素贝叶斯}（naive Bayes）采用生成式路线，按类别估计各特征的分布。它假设给定
类别后，知道一个特征的值不会进一步改变其他特征的分布，也就是特征在该条件下相互独立。
这样就能用各特征的概率相乘来简化联合概率。\term{高斯判别分析}
（Gaussian discriminant analysis）也采用生成式路线，但用多元高斯分布描述各类特征，
从而保留特征之间的相关性。“生成式”的名称来自模型所描述的数据产生过程：
先按$p(y)$选择类别，再按$p(\bm{x}\mid y)$生成特征。在本章的邮件表示下，
生成的是促销词和链接的计数；若要生成完整邮件文本，还需要对文本本身建立分布。

图\ref{fig:models-discriminative-generative}比较了两条路线。它们都能给出类别后验，
区别在于这个后验是直接学习的，还是由所学习的联合分布推导出来的。

\begin{figure}[htbp]
  \centering
  \begin{tikzpicture}[x=1mm,y=1mm,font=\footnotesize,
  cell/.style={book node,font=\footnotesize,minimum height=10mm,inner sep=5pt}]
  \node[anchor=west,font=\heiti\small] at (0,30) {直接学习类别条件分布};
  \node[cell] (a) at (24,17) {$p(y\mid\bm{x})$};
  \node[cell,fill=FigureOutput!10] (b) at (83,17) {类别概率};
  \draw[book arrow] (a)--node[above] {直接计算}(b);
  \node[anchor=west,font=\heiti\small] at (0,-2) {先描述联合分布};
  \node[cell] (c) at (24,-16) {$p(y)\,p(\bm{x}\mid y)$};
  \node[cell,fill=FigureOutput!10] (d) at (83,-16) {类别概率};
  \draw[book arrow] (c)--node[above] {贝叶斯公式}(d);
\end{tikzpicture}

  \caption{同一分类任务的两条概率建模路线。判别式模型直接学习条件类别分布；
  生成式模型先学习类别比例与类内特征分布，再通过贝叶斯公式得到类别后验。}
  \label{fig:models-discriminative-generative}
\end{figure}

两条路线的差异也体现在损失与学习准则中。假设$n$个带标签样本$(\bm{x}_i,y_i)$相互独立，
来自同一分布，用$\bm{\theta}$分别统称各模型自己的参数。
训练时可以比较不同参数给已观测数据分配了多大的概率或密度；把数据固定、把这一数值
看成参数的函数，就得到\term{似然}（likelihood）。独立样本的似然相乘，取对数后则
变成相加；再加上负号，最大化似然就变成最小化负对数似然。判别式与生成式模型采用
这条准则时，分别评价条件分布与联合分布：
\begin{align}
  \ell_{\mathrm{disc}}(\bm{\theta})
    &=-\frac{1}{n}\sum_{i=1}^{n}
      \log p_{\bm{\theta}}(y_i\mid\bm{x}_i),
    \label{eq:models-discriminative-objective}\\
  \ell_{\mathrm{gen}}(\bm{\theta})
    &=-\frac{1}{n}\sum_{i=1}^{n}
      \log p_{\bm{\theta}}(\bm{x}_i,y_i).
    \label{eq:models-generative-objective}
\end{align}
前者只评价给定输入后，模型为实际类别分配了多少概率；后者还评价输入与类别共同出现的
可能程度。因此，即使两种模型最终给出的类别后验具有相同的函数形式，学习出的参数也不一定相同。

当任务只要求分类时，判别式模型可以集中学习与类别有关的关系，减少对输入分布形状的
额外假设。生成式模型描述的范围更广，还可以在相应分布假设下进行采样或补全缺失特征；
代价是需要估计更多分布关系。额外假设可能帮助有限数据下的估计，也可能因与数据不符而
影响分类。例如，朴素贝叶斯的条件独立假设简化了估计，但相关特征可能使其概率估计偏离实际。
两条路线的表现取决于具体模型、数据与估计方法，逻辑回归和朴素贝叶斯的比较提供了一个
研究实例\scite{ng2001}

这两条路线都可以得到类别概率，区别在于训练时需要描述多大范围的分布关系。下面暂且
认为这些分布的参数已经确定，看看如何从中计算任务需要的结果；参数本身的不确定性
将在随后讨论。

\paragraph{从函数计算到概率推断}

无论类别后验是直接建立的，还是由联合分布导出的，使用概率模型时都需要从已建立的分布
计算当前任务关心的量。这个过程称为\term{概率推断}
（probabilistic inference）。例如，在温度预测中，要求温度落在某个区间的概率，需要将
该区间内的概率密度积分：
\begin{equation}
  \mathbb{P}(20\leq y\leq24\mid\bm{x})
  =\int_{20}^{24}p(y\mid\bm{x})\,\mathrm{d}y.
  \label{eq:models-predictive-interval}
\end{equation}
这里温度以摄氏度计。在聚类中，GMM先描述各组内的数据分布，再根据已经观察到的样本
计算其组别概率。这两种计算都从已有的概率关系出发，得到当前问题需要的分布或概率值。

温度的点预测只需代入函数求值，区间概率却需要把区间内的可能结果累积起来。更一般地，
概率推断可能需要组合事件的概率，对尚未观测的变量求和或积分，以及根据新观测计算条件
分布。函数仍然可以计算均值等参数，任务所要求的概率则由相应分布进一步得到。

当许多变量相互依赖时，概率推断的计算量可能迅速增加。例如，若必须逐一考虑
$m$个二值变量的取值组合，就有$2^m$种情况；连续变量则可能涉及高维积分。
训练还可能反复进行这些计算，才能评价模型或更新参数。因此，复杂概率模型的训练与推断
常需借助采样或较简单的分布进行近似。
简单分布则可能直接计算，例如上述高斯预测只需均值与方差，就能确定整条密度曲线。

当变量较多时，还需要一种方式组织这些依赖关系。\term{概率图模型}
（probabilistic graphical model）用图表示多个变量的概率关系：
节点代表随机变量，连线及其缺失表达分布如何分解、哪些变量在给定条件下相互独立。
这里的条件独立表示：给定某些变量后，得知另一个变量的值，不再改变所考察变量的条件分布。
这种结构有助于分解计算，不过复杂图中的推断仍可能需要近似方法。
概率图的具体定义与计算将在第\ref{part:probabilistic-models}部分展开\scite{goodfellow2016}

\paragraph{参数的点估计与贝叶斯推断}

前一小节讨论的是模型参数给定以后怎样进行概率计算。回到训练过程，还需要决定怎样根据
数据确定这些参数。
一种做法是选择一组具体数值，例如估计温度模型的权重与噪声方差，然后用这组数值预测。
用一个值表示每个未知参数，称为\term{点估计}（point estimation）。前面介绍的似然
准则可以用于这一选择：在允许的参数范围内寻找使似然最大的取值，称为
\term{最大似然估计}（maximum likelihood estimation）。参数取确定值之后，温度仍然
由一个分布描述；确定的是分布的参数，并非明天的温度。

如果训练数据很少，不同参数都可能合理地解释已有观测。这时可以进一步用分布描述参数
的不确定性。\term{贝叶斯推断}（Bayesian inference）用数据更新关于未知量的概率描述。
对于模型参数，\term{先验分布}（prior distribution）$p(\bm{\theta})$表达观察本批数据
之前的假设；结合数据$\mathcal{D}$之后，得到\term{后验分布}（posterior distribution）
$p(\bm{\theta}\mid\mathcal{D})$。预测时，不只采用一组参数，而是将不同参数给出的
结果按后验权重综合起来：
\begin{equation}
  p(y\mid\bm{x},\mathcal{D})
  =\int p(y\mid\bm{x},\bm{\theta})
    p(\bm{\theta}\mid\mathcal{D})\,\mathrm{d}\bm{\theta}.
  \label{eq:models-bayesian-prediction}
\end{equation}
在温度例子中，固定参数下的分布描述相似气象条件仍可能对应不同温度；参数后验还反映了
有限数据使我们难以确定预测关系。这两个来源都可以影响最终的预测区间。
贝叶斯推断将概率推断用于参数，代价是通常需要额外计算或近似参数后验。

同一个概率模型因此有两种参数处理方式：用点估计固定分布，或用参数后验综合不同分布。
概率图中的变量有概率描述，也不要求所有分布参数都采用后一种处理。
两种选择还可以与不同容量结合：固定成分数的GMM是参数化概率模型；\term{高斯过程}
（Gaussian process）对可能的函数建立概率描述，不预先把所有函数限制在固定数量的系数中，
是非参数化概率模型的例子。非参数化所讨论的表示规模，与是否建立概率分布可以分别选择\scite{rasmussen2006}

\paragraph{概率描述带来了什么}

建立分布与进行推断增加了计算，得到的概率信息能帮我们作什么决定？例如，两个模型都
预测明天$22\,^{\circ}\mathrm{C}$，
但一个把大部分概率集中在$21$到$23\,^{\circ}\mathrm{C}$，另一个分布在更宽范围内。
安排对温度敏感的活动时，这两种预测提供的信息不同。GMM给出的组别概率也能帮助发现
难以明确归类的样本。

这些信息是否可靠，要用未参与训练的数据检验。若模型反复给出名义覆盖概率为$90\%$的
预测区间，就应检查实际温度是否在大约$90\%$的情况下落入相应区间。在覆盖率相当的
前提下，区间更窄通常能提供更具体的判断。始终给出很宽的区间虽然容易覆盖真实值，
却提供不了多少判断依据。
因此，评价概率模型既要看中心预测是否准确，也要看不确定性的大小是否与实际误差相称。

分布假设也会限制效果。例如，若同类条件下的温度存在两个明显峰值，单个高斯分布可能把
均值放在两个峰之间，并给出不合适的区间。聚类时，错误的群组形状假设也会影响分配概率。
概率模型增加了描述和使用不确定性的能力，是否改善具体任务，仍取决于这些假设与数据
是否吻合，以及分布能否得到可靠估计。

\section{变量的关系形式}
\label{sec:models-structure}

给定一个输入，模型究竟通过什么关系算出结果？即使都输出温度，预测也可以随着某个特征
按固定幅度变化，或在不同范围内采用不同的变化规律。这是模型结构关心的问题。
数值预测可以考察预测函数，分类可以考察得分或决策边界，降维可以考察输入到低维坐标
的映射。判断线性与非线性之前，需要明确正在看哪一条关系。

\subsection{线性结构}
\label{sec:models-linear}

一种直接的假设是：每个特征的作用可以用固定权重表示，再把各项贡献相加。线性回归
正是将$d$维输入$\bm{x}$按权重$\bm{w}\in\mathbb{R}^{d}$加权求和，加上偏置$b$，
得到预测值：
\begin{equation}
  \hat{y}=\bm{w}^{\top}\bm{x}+b
  =\sum_{j=1}^{d}w_jx_j+b.
  \label{eq:models-linear-score}
\end{equation}
固定其他特征后，$x_j$每增加一个单位，预测值都改变$w_j$。
例如，用面积预测房价时，线性关系假设每增加一平方米带来的预测增量相同。
严格说，带常数项$b$的关系称为仿射关系；机器学习中通常也将它归入线性模型。

概率需要落在零与一之间，任意的加权和却没有这个限制。逻辑回归在加权和之后加上一个
转换函数。对于二分类，记正类概率为$p$，则
\begin{equation}
  p=\frac{1}{1+\exp[-(\bm{w}^{\top}\bm{x}+b)]},
  \qquad
  \log\frac{p}{1-p}=\bm{w}^{\top}\bm{x}+b.
  \label{eq:models-logistic-structure}
\end{equation}
左式将任意实数得分转为零与一之间的概率。右式来自对左式的变形：$p/(1-p)$比较正类
与负类的概率大小，再取对数，就得到\term{对数几率}（log odds）。这个量仍然是输入
的仿射函数。用$p=0.5$作分类阈值时，两类概率相等，对数几率为零，边界因而满足
$\bm{w}^{\top}\bm{x}+b=0$，在二维平面中就是一条直线。
逻辑回归因此属于线性分类模型，尽管其输出概率包含非线性函数。

这里的线性描述的是对数几率与边界，在线性回归中描述的则是预测值。基础
\term{感知机}（perceptron）也用加权和作为分类得分，按得分的正负判断两类。
它与线性SVM的训练准则不同，却都用一个线性分界面完成二分类。

线性结构也可以用于表示学习。PCA的线性体现在投影上：固定均值与投影方向后，每个低维
坐标都是输入减去均值后的加权和。
PPCA增加了分布假设，但从低维表示到观测均值的关系仍然是线性的。
因此，线性结构既能用于数值预测和分类，也能用于降维与概率建模。

聚类没有预设的类别边界，但把样本分给最近中心时也会产生分界。对于K-means的两个中心
$\bm{\mu}_1$与$\bm{\mu}_2$，边界上的样本满足
$\lVert\bm{x}-\bm{\mu}_1\rVert_2^2=\lVert\bm{x}-\bm{\mu}_2\rVert_2^2$。
展开两边并消去共同的$\bm{x}^{\top}\bm{x}$项，得到
\begin{equation}
  2(\bm{\mu}_2-\bm{\mu}_1)^{\top}\bm{x}
  =\lVert\bm{\mu}_2\rVert_2^2-\lVert\bm{\mu}_1\rVert_2^2.
  \label{eq:models-kmeans-boundary}
\end{equation}
这是关于输入的仿射方程，所以标准K-means任意两个不同中心之间的分配边界是线性的。
但样本到中心的平方距离是二次函数，最终的簇编号也不是输入的线性函数。
后文将K-means放在线性一侧时，专指这种两两分配边界。

判断线性结构时，还要明确模型相对于哪一组输入特征是线性的。若先把$x$变成$x^2$，
再计算$wx^2+b$，模型对新特征$x^2$是线性的，对原始输入$x$则是非线性的。
本章以原始输入中的关系为主要依据。“参数以线性方式出现”与“预测对输入是线性的”
需要结合具体公式区分。

\subsection{非线性结构}
\label{sec:models-nonlinear}

上一节的模型用一个仿射关系连接输入与所考察的量，但现实中的许多关系会随条件改变。
车辆速度增加时，刹车距离通常不会按固定增量增长；
同一种药物增加剂量，在不同剂量区间可能产生不同幅度的作用；邮件中的促销词是否可疑，
也可能取决于同时出现了多少外部链接。这些关系需要表达弯曲、阈值或特征之间的共同作用。

沿用上一节的判断依据，如果所考察的预测值或得分无法写成原始输入的仿射函数，
就属于非线性关系。对于二分类，无法用单一超平面表示的弯曲或分段边界，也需要非线性结构。
图\ref{fig:models-linear-nonlinear}展示了线性与非线性关系在回归和分类中的区别。

\begin{figure}[!ht]
  \centering
  % 原创教学几何示意；四幅图在正文宽度内按两行两列排列，固定坐标与解析曲线。
\begin{tikzpicture}[x=1mm,y=1mm,font=\footnotesize]
  \foreach \xshift/\yshift/\title in {
    0/52/线性预测,55/52/非线性预测,
    0/0/线性分类,55/0/非线性分类}{
    \begin{scope}[xshift=\xshift mm,yshift=\yshift mm]
      \node[font=\heiti\small] at (15,37) {\title};
      \fill[BookPaper] (0,0) rectangle (30,28);
    \end{scope}
  }
  \foreach \xshift in {0,55}{
    \begin{scope}[xshift=\xshift mm,yshift=52mm]
      \draw[book arrow] (0,0)--(33,0) node[right] {$x$};
      \draw[book arrow] (0,0)--(0,30) node[above] {$\hat y$};
    \end{scope}
  }
  \foreach \xshift in {0,55}{
    \begin{scope}[xshift=\xshift mm]
      \draw[book arrow] (0,0)--(33,0) node[right] {$x_1$};
      \draw[book arrow] (0,0)--(0,30) node[above] {$x_2$};
    \end{scope}
  }
  \begin{scope}[yshift=52mm]
    \draw[FigureModel,very thick] (2,4)--(28,25);
  \end{scope}
  \begin{scope}[xshift=55mm,yshift=52mm]
    \draw[FigureModel,very thick]
      plot[domain=2:28,samples=60] ({\x},{4+0.032*(\x-2)^2});
  \end{scope}
  \begin{scope}
    \draw[FigureModel,thick] (2,4)--(28,24);
    \foreach \x/\y in {4/13,8/20,13/23,19/26}{
      \fill[FigureInput] (\x,\y) circle (1);
    }
    \foreach \x/\y in {9/4,17/7,24/10,28/17}{
      \fill[FigureLoss] (\x-0.9,\y-0.9) rectangle (\x+0.9,\y+0.9);
    }
  \end{scope}
  \begin{scope}[xshift=55mm]
    \draw[FigureModel,thick] (15,14) circle (8);
    \foreach \x/\y in {11/12,16/18,19/12,14/8}{
      \fill[FigureInput] (\x,\y) circle (1);
    }
    \foreach \x/\y in {4/6,5/25,16/26,28/22,27/5}{
      \fill[FigureLoss] (\x-0.9,\y-0.9) rectangle (\x+0.9,\y+0.9);
    }
  \end{scope}
  \node at (15,44) {固定增量};
  \node at (70,44) {增量随输入改变};
  \node at (15,-8) {直线分界};
  \node at (70,-8) {曲线分界};
\end{tikzpicture}

  \caption{线性与非线性关系的几何示意。上排依次为线性预测与非线性预测，
  下排依次为线性分类与非线性分类；圆点和方点代表两类。非线性描述函数或边界的形状，
  两类模型都可以定义在同一个输入空间中。}
  \label{fig:models-linear-nonlinear}
\end{figure}

图中的弯曲关系可以由不同模型得到。固定次数的多项式回归对原始输入是非线性的，但只需
固定数量的系数，属于参数化模型。高斯过程回归同样可以表达弯曲关系，却不受固定有限维
系数限制，还为函数值建立概率分布。曲线的形状说明了关系形式；它本身没有告诉我们
参数化方式与表示规模、概率语义如何选择。

非线性也不只表现为单个变量对应一条弯曲曲线。若模型包含$x_1x_2$，那么$x_1$对输出的
影响会随$x_2$而改变，这称为特征之间的交互。回到邮件例子，少量促销词本身可能并不可疑，
但当外部链接很多时，同样数量的促销词可能明显提高垃圾邮件得分。单独查看两个特征时，
这种共同作用不一定明显；把它们放在同一关系中，才会看到线性加权和无法表达的条件变化。

是否需要非线性结构，应结合任务和数据判断。线性模型提供了清楚的基准：每个特征通过固定
权重影响所考察的量。若验证数据表明误差仍随输入区域呈现稳定的弯曲、阈值或交互模式，
再引入分段规则或非线性变换会更有依据。更灵活的结构扩大了可表达的关系，也需要结合
数据量与结构约束，防止把偶然波动当成规律。常见的做法是按条件划分区域，或为输入
增加非线性变换。
\FloatBarrier

\paragraph{按条件采用不同规则}

表达非线性的一种直接办法，是先划分输入空间，再在不同区域使用不同规则。决策树通过
一系列条件判断，将不同输入交给不同规则处理。
例如，先检查促销词数量；只有数量达到阈值时，再检查外部链接数。
一条从根到叶的路径对应一组条件，不同叶节点对输入空间的不同区域给出预测。
图\ref{fig:models-tree-partition}将同一组规则画成树与平面分区。

\begin{figure}[htbp]
  \centering
  \begin{tikzpicture}[x=1mm,y=1mm,font=\footnotesize,
  rule/.style={fill=BookPaper,inner sep=4pt,align=center}]
  \node[rule] (r) at (22,37) {$x_1\geq2$？};
  \node[rule] (a) at (9,20) {类别 $0$};
  \node[rule] (b) at (35,20) {$x_2\geq1$？};
  \node[rule] (c) at (24,3) {类别 $0$};
  \node[rule,fill=FigureModel!12] (d) at (46,3) {类别 $1$};
  \draw[book arrow] (r)--node[left] {否}(a);
  \draw[book arrow] (r)--node[right] {是}(b);
  \draw[book arrow] (b)--node[left] {否}(c);
  \draw[book arrow] (b)--node[right] {是}(d);
  \begin{scope}[xshift=64mm]
    \fill[BookPaper] (0,0) rectangle (37,36);
    \fill[FigureModel!18] (18,12) rectangle (37,36);
    \draw[book arrow] (0,0)--(40,0) node[right] {$x_1$};
    \draw[book arrow] (0,0)--(0,40) node[above] {$x_2$};
    \draw[BookInk,thick] (18,0)--(18,36);
    \draw[BookInk,thick] (18,12)--(37,12);
    \node[below] at (18,0) {$2$};
    \node[left] at (0,12) {$1$};
    \node at (8,23) {$0$};
    \node at (28,6) {$0$};
    \node at (28,25) {$1$};
  \end{scope}
\end{tikzpicture}

  \caption{按条件分情况处理输入。左侧为判断树，右侧为对应的输入区域。
  该教学例子在$x_1\geq2$且$x_2\geq1$时输出类别$1$，其余区域输出$0$。}
  \label{fig:models-tree-partition}
\end{figure}

常见决策树每次只比较一个特征与一个阈值，在叶节点输出固定值。
单次比较虽然简单，多次分支组合后却能形成分段的预测关系和复杂边界。
这种方式便于表达“在某种条件下采用某条规则”，但相邻区域的预测也可能出现跳变。
若允许叶节点数随训练数据增加，树属于非参数化设定；若预先限制叶节点数量，表示规模
也就受到固定上限约束。这是结构与容量需要分别考察的一个例子。

\paragraph{对输入进行非线性变换}

除了显式划分区域，还可以先变换输入，再计算输出。例如，加入$x_1^2$可以表达弯曲关系，
加入$x_1x_2$可以表达两个特征的共同作用。用$\bm{\phi}(\bm{x})$表示变换后的特征，得分可写为
\begin{equation}
  s(\bm{x})=\bm{w}^{\top}\bm{\phi}(\bm{x})+b.
  \label{eq:models-feature-score}
\end{equation}
只要变换包含非线性关系，线性得分就可能在原始空间中对应非线性边界。
例如，中心区域是一类、外围是一类时，可以把到原点的平方距离
$x_1^2+x_2^2$作为新特征，再用一个阈值区分两类。

\term{核方法}（kernel methods）也使用这种思路：通过计算样本在特征空间中的内积，
间接使用变换后的关系，无需逐项列出所有新特征。线性核仍对应线性结构，高斯核等非线性核
则能够形成弯曲边界。
标准高斯核SVM可以随样本增加保留更多表示系数，属于非参数化设定；
若预先选取固定数量的变换特征进行近似，则成为另一种容量设定\scite{cortes1995}

手工选择$x^2$或$x_1x_2$，要求我们预先猜测哪些变换有用。神经网络让中间变换也参与
学习。与基础感知机只有一个线性得分相比，多层网络在中间层加入非线性函数，通常称为
\term{激活函数}（activation function）；例如，把负值置零、保留非负值，就能让不同
输入范围采用不同的计算关系。
网络将多次变换组合起来，根据训练目标共同调整中间表示和最终预测。
如果各层都只有加权求和而没有非线性操作，把前一层代入后一层以后，仍能合并成一次
加权求和与偏置。因此，仅仅增加这类层数并不能得到非线性关系。

树直接划分区域，便于表达条件规则；特征变换则改变输入表示，使弯曲关系或共同作用可以
通过新的特征表达。两种思路也会交叉，例如神经网络的分段线性激活同样把输入划分成不同
计算区域。能否把这些表达能力用于新样本，仍取决于数据、结构限制与训练方法。

神经网络与前文的概率图模型都可以画成图，但图中的连线承担不同作用：网络连线说明
数值如何参与下一步计算，概率图的连线及其缺失说明变量间的概率依赖。一个概率模型可以
用神经网络计算分布参数；此时非线性计算结构与概率描述就在同一个模型中结合起来。

\FloatBarrier
\section{输出的形式与含义}
\label{sec:models-output}

模型内部采用的关系已经确定，交给使用者的结果却仍有选择。温度模型可以返回一个数，
也可以返回一个区间或完整分布；分类器可以返回类别，也可以保留各类概率。
理解输出时需要回答两个不同问题：结果怎样组织，以及这些数值代表什么。前者区分标量、
向量和带有位置或顺序的结构化对象；后者区分点预测、类别编码、分布摘要、完整分布和
采样结果。只看其中一项，不能确定模型向使用者提供了哪些信息。

\subsection{输出的组织形式}
\label{sec:models-output-form}

\paragraph{标量输出}
\label{sec:models-scalar-output}

用一个数表示结果，便是\term{标量输出}（scalar output）。温度预测、价格预测和某个
商品的评分都可以采用这种形式。
这类数值具有对应的单位或尺度，例如$22\,^{\circ}\mathrm{C}$可以直接用于温度比较。

同样的标量形式也可以编码离散结果。邮件分类可以用$0$表示正常、$1$表示垃圾邮件；
聚类可以用$1,2,\ldots,k$区分不同群组。分类编号的语义由标签定义，聚类编号则标记
从数据中发现的群组。交换两个簇的编号不会改变分组结果，因此编号之间的数值差通常
没有距离意义。

\paragraph{向量与结构化输出}
\label{sec:models-vector-output}

当一个结果需要同时保留多个分量时，可以把它们按顺序组成\term{向量输出}
（vector output）。例如，同时预测温度、湿度和风速，得到三个数值；预测未来七天的
温度，得到按日期排列的七个数值。各分量对应不同目标，
使用时需要知道它们的顺序、单位和含义。

另一类向量用来表示输入本身。PCA输出的每个坐标表示样本沿一个投影方向的位置；
神经网络也可以把一张图像变成一个特征向量，并通过适当的训练目标让内容相近的图像
具有相近的表示。
这些向量随后可以用于检索、分类或聚类，而不必直接对应某个最终答案。

图像重构也可以看作输出一组数值，只是它们按像素位置和颜色通道排列成数组。
文本输出则还带有符号顺序与长度等组织信息。讨论这类结果时，除了分量的取值，
还需要保留它们之间的位置、顺序或连接关系。

标量、向量和结构化对象只说明结果怎样组织，不说明其中数值的含义。长度为三的向量
可能表示三个温度预测，也可能表示三个类别的概率；后者还要求各分量非负、总和为一。
同样，一个标量既可能是温度预测，也可能只是类别或簇的编号。要继续理解输出，还要
说明这些数值代表一个直接结果、对分布的概括，还是完整的概率描述。

\subsection{输出的含义}
\label{sec:models-output-semantics}

第\ref{sec:models-probability}节考察模型是否建立概率分布；这里进一步考察接口是否把
完整分布交给使用者，还是只返回一个结果或摘要。模型内部的概率语义与接口交付的信息
有关，却不是同一个判断。

\paragraph{点结果与分布摘要}

模型可以直接返回一个点预测、类别或簇编号。概率模型也可以只交付一个便于使用的摘要：
逻辑回归能够返回概率最大的类别，高斯温度模型能够返回分布均值或预测区间。
这些输出分别可能采用标量、向量或区间等形式，但都没有把模型保留的全部概率信息交给
使用者。摘要便于比较和决策，代价是其他可能结果及其概率不再完整可见。

\paragraph{概率分布与采样}
\label{sec:models-distribution-output}

如果需要知道各个可能结果分别有多大概率，模型就要提供\term{概率分布输出}
（probability distribution output），保留各种结果的可能程度。对于有限类别，
可以逐项给出非负、总和为一的概率。
例如，邮件分类输出$(0.2,0.8)^{\top}$，表示正常邮件与垃圾邮件的概率分别为$0.2$和$0.8$。
相比只返回“垃圾邮件”，这一结果还体现了模型对判断的确定程度。

连续变量有无穷多个可能取值，通常用分布的形式和参数来表示。
例如，温度模型指定高斯分布，并输出均值$22$和方差$4$，就确定了相应的密度曲线。
密度曲线下某个区间的面积给出温度落入该区间的概率；曲线在某一点的高度表示密度，
并不是温度恰好等于该数值的概率。

同一个概率模型可以提供不同接口：直接返回预测结果的概率分布，计算前述摘要，
或者从分布中采样。对于温度预测，摘要可以是均值$22\,^{\circ}\mathrm{C}$，也可以是具有
指定覆盖概率的预测区间；采样则可能得到$21.3\,^{\circ}\mathrm{C}$等一个可能的温度。
对于图像生成，模型的一次采样得到一张图像，多次采样可以得到不同图像。
分布保留各种可能结果，摘要便于概括和决策，采样则用于展示或模拟具体情形。

输出的组织形式与结果含义因此是两条不同的描述轴。类别分布用向量表示，分布均值用
标量表示，一次采样则可能产生高维数组。说明返回了几个数之后，还要说明这些数代表
点预测、分布摘要、完整分布，还是从分布中取出的一个样本。

\section{损失与学习准则}
\label{sec:models-objectives}

知道模型要输出什么，还不足以开始训练：我们需要约定怎样判断某个结果更好。
例如，分类与聚类都可以把邮件特征映射为一个编号。分类已有真实标签，希望预测类别
与标签一致；聚类没有这样的标签，需要根据样本之间的关系形成分组。
两者返回的都是编号，学习时评价结果的依据却不同。

本节按评价所依赖的信息组织三类问题。有直接目标时，可以逐个样本比较模型输出与目标，
再根据输出对象选择数值、向量或概率分布的差异。若没有逐样本答案，或者任务还关心
多个样本与表示之间的关系，就需要规定紧凑性、间隔、相对邻近或低维结构。当前行动
还会影响后续结果时，评价范围则要扩展到整段轨迹和长期收益。
这些准则可以合用，例如PCA既限制表示所在空间的维数，又衡量重构向量与原始输入的误差。

\subsection{逐样本的目标差异}
\label{sec:models-pointwise-objectives}

先看单个结果有已知答案的情形。温度、价格和销量预测都有一个可以比较的目标值；
分类也有待预测的真实类别。
这类任务希望模型的结果接近已知答案，但连续数值与离散类别需要不同的比较方式。
记第$i$个样本的真实值为$y_i$、预测为$\hat{y}_i$，通常先计算每个样本的损失，
再取$n$个样本上的平均值。

\paragraph{连续数值}

预测温度偏高两度与偏低两度，都有两个单位的偏差。为了避免正负误差在求和时互相抵消，
可以用平方损失或绝对损失评价单个预测：
\[
  \ell_{\mathrm{sq}}(\hat y,y)=(\hat y-y)^2,\qquad
  \ell_{\mathrm{abs}}(\hat y,y)=|\hat y-y|.
\]
分别在样本上取平均，就得到第~\ref{sec:performance-measures}节介绍的均方误差和平均绝对误差。
若实际温度为$20$、预测为$22$，单个样本的平方误差为$4$，绝对误差为$2$。
当预测偏差由$2$增至$4$，平方误差变成原来的四倍，绝对误差则变成两倍。
平方误差更重视较大的偏差，且作为预测值的函数处处可导；绝对误差对少数大偏差更稳健，
但在误差为零处不可导，优化时需要相应处理。

如果真实错误的代价不对称，还可以分别设置高估和低估的损失。
例如，备货不足与库存过多的成本不同，选择误差形式时就应体现这种差别。
因此，“接近真实值”还需要进一步说明：哪些方向、多大程度的偏离更值得惩罚。

\paragraph{离散类别}

离散数值通常用于编码类别，而类别编号只表示身份。若$1$、$2$、$3$分别代表猫、狗和鸟，
编号之差不能衡量分类错误的严重程度。最直接的\term{零一损失}（zero-one loss）只检查
预测是否正确：
\begin{equation}
  \ell_{01}(\hat{y},y)=
  \begin{cases}
    0, & \hat{y}=y,\\
    1, & \hat{y}\ne y.
  \end{cases}
  \label{eq:models-zero-one-loss}
\end{equation}
它在数据集上的平均值就是分类错误率。不过，只要预测类别没有改变，参数调整得再多，
零一损失也保持原值，因而难以指引连续参数怎样更新。训练时常改用随得分或概率变化的
\term{代理损失}（surrogate loss），用更便于优化的目标帮助减少分类错误。常见选择包括
后文介绍的类别概率交叉熵，以及第\ref{sec:models-spatial-objectives}节的分类得分间隔损失。
如果不同类型的误判具有不同代价，也可以用事先给定的代价表替代统一的$1$。

\paragraph{向量的差异}
\label{sec:models-vector-objectives}

当目标由多个分量组成时，需要比较两个向量而不是两个标量。图像去噪希望输出接近干净
图像，压缩与重构希望恢复的内容接近原始输入，多变量预测希望
一组预测共同接近真实观测。这些任务都有两个需要比较的向量。
表示学习也会比较向量：例如，让同一图像的不同视图得到相近的特征表示。
前者通常关心每个分量恢复得怎样，后者则可能更关心整体方向或相对位置。

对于预测向量$\hat{\bm{y}}$与目标向量$\bm{y}\in\mathbb{R}^{m}$，常用的两种误差为
\begin{align}
  \ell_2(\hat{\bm{y}},\bm{y})
    &=\lVert\hat{\bm{y}}-\bm{y}\rVert_2^2
      =\sum_{j=1}^{m}(\hat{y}_j-y_j)^2,
    \label{eq:models-vector-loss}\\
  \ell_1(\hat{\bm{y}},\bm{y})
    &=\lVert\hat{\bm{y}}-\bm{y}\rVert_1
      =\sum_{j=1}^{m}|\hat{y}_j-y_j|.
    \label{eq:models-vector-l1}
\end{align}
平方欧氏误差将各分量的平方偏差相加，较大偏差占更高权重；$L_1$误差累计绝对偏差，
较少被个别异常分量主导。两者都依赖分量的尺度，例如同时比较温度与气压时，需要考虑
单位差异，决定是否先做缩放或设置权重。

若主要关心方向，可以使用\term{余弦相似度}（cosine similarity），即两个向量夹角的
余弦值：方向越接近，这个值越大。
对于非零向量$\bm{a}$与$\bm{b}$，相应的损失为
\begin{equation}
  \ell_{\mathrm{cos}}(\bm{a},\bm{b})
  =1-\frac{\bm{a}^{\top}\bm{b}}
          {\lVert\bm{a}\rVert_2\lVert\bm{b}\rVert_2}.
  \label{eq:models-cosine}
\end{equation}
同向向量的损失为零。例如，$(1,2)^{\top}$与$(2,4)^{\top}$长度不同，但方向相同，
所以余弦损失为零，欧氏误差却不为零。图\ref{fig:models-vector-comparison}展示了这一差别。
比较文档或图像的特征表示时，这种对长度变化不敏感的度量可能有用；
若向量的幅度表示真实温度、亮度或数量，忽略长度就可能丢失任务所需的信息。

\begin{figure}[htbp]
  \centering
  % 原创教学几何示意：同向不同长度与等长不同方向。
\begin{tikzpicture}[x=1mm,y=1mm,font=\heiti\small]
  \node at (21,40) {比较端点距离};
  \node at (79,40) {比较方向};
  \foreach \shift in {0,58}{
    \begin{scope}[xshift=\shift mm]
      \draw[book arrow] (0,0)--(40,0) node[right] {$x_1$};
      \draw[book arrow] (0,0)--(0,33) node[above] {$x_2$};
    \end{scope}
  }
  \draw[book arrow,FigureModel,thick] (0,0)--(21,28) node[above right] {$\bm b$};
  \draw[book arrow,FigureInput,thick] (0,0)--(10.5,14) node[left] {$\bm a$};
  \draw[FigureLoss,thick,dashed] (13.5,13)--(24,27);
  \draw[FigureLoss] (12.2,14)--(14.8,12) (22.7,28)--(25.3,26);
  \node[text=FigureLoss,align=center] at (32,16) {长度\\不同};
  \begin{scope}[xshift=58mm]
    \draw[book arrow,FigureInput,thick] (0,0)--(28,7) node[right] {$\bm a$};
    \draw[book arrow,FigureModel,thick] (0,0)--(14,25.24) node[above right] {$\bm b$};
    \draw[FigureLoss,thick] (9.70,2.43) arc[start angle=14.04,end angle=60.98,radius=10];
    \node[text=FigureLoss] at (14,10) {$\alpha$};
  \end{scope}
  \node at (21,-9) {同向也可以相距很远};
  \node at (79,-9) {夹角越小，方向越接近};
\end{tikzpicture}

  \caption{两种向量比较方式关注不同差异。左侧比较两个端点之间的距离，
  右侧比较向量方向之间的夹角。同向但长度不同的向量，余弦损失为零，欧氏误差仍大于零。}
  \label{fig:models-vector-comparison}
\end{figure}

还需要判断数值误差能否代表任务中的相似程度。
例如，一张图像仅平移一个像素，内容可能几乎没变，却会造成许多像素值不一致。
直接比较像素适合要求位置精确对应的重构；若关心物体内容，可以先提取相应特征，
再比较特征向量。度量与所比较的表示需要一起选择。

\paragraph{概率分布的差异}
\label{sec:models-probability-objectives}

如果模型输出概率分布，只比较最终类别或分布均值会丢失其余概率信息。邮件分类可能不仅
要求类别正确，还要求预测概率可信；学习数据分布时，希望模型描述的
样本出现规律接近实际数据；让一个模型模仿另一个模型时，也可能希望两者给出的类别
概率接近。这些任务比较的是概率分布，而不只是最终选出的一个答案。

先考虑有限类别。设目标分布为$\bm{p}=(p_1,\ldots,p_k)^{\top}$，预测分布为
$\bm{q}=(q_1,\ldots,q_k)^{\top}$，二者的分量都非负且分别求和为一。
模型给某一结果的概率越低，$-\log q_j$就越大；再按目标分布中该结果的概率$p_j$
加权，就得到\term{交叉熵}（cross-entropy）。它衡量模型对目标分布下的结果平均赋予了
多大的负对数概率：
\begin{equation}
  \operatorname{CE}(\bm{p},\bm{q})=-\sum_{j=1}^{k}p_j\log q_j.
  \label{eq:models-cross-entropy}
\end{equation}
如果已知样本属于类别$c$，可以令$p_c=1$、其余分量为零；这种编码称为
\term{独热表示}（one-hot representation）。此时交叉熵就是$-\log q_c$。
给正确类别$0.9$的概率比给$0.1$的损失更小；当$q_c$趋近零时，损失趋向无穷大，
因而非常自信却判错的预测会受到强烈惩罚。这里的对数采用自然底数。

交叉熵在预测分布等于目标分布时，也未必为零。如果希望用零表示两个分布完全一致，
可以减去目标分布与自身的交叉熵。

两者相减得到\term{KL散度}（Kullback--Leibler divergence），用来衡量两个分布的差异：
\begin{equation}
  D_{\mathrm{KL}}(\bm{p}\Vert\bm{q})
  =\sum_{j=1}^{k}p_j\log\frac{p_j}{q_j}.
  \label{eq:models-kl}
\end{equation}
约定$p_j=0$的项为零；若$p_j>0$而$q_j=0$，散度为无穷大。
它非负，并在两个分布相同时为零。把式\eqref{eq:models-kl}中的对数比值拆开，就得到
$D_{\mathrm{KL}}(\bm{p}\Vert\bm{q})
=\operatorname{CE}(\bm{p},\bm{q})-\operatorname{CE}(\bm{p},\bm{p})$。
固定目标$\bm{p}$时，第二项不随预测$\bm{q}$变化，所以最小化KL散度与交叉熵会得到
相同的优化要求。
KL散度通常不对称：交换目标与预测，关注的概率区域和得到的数值都可能改变。

也可以直接比较概率向量的平方差。对于独热目标，这就是\term{Brier评分}（Brier score）：
\begin{equation}
  \ell_{\mathrm{Brier}}(\bm{p},\bm{q})
  =\sum_{j=1}^{k}(q_j-p_j)^2.
  \label{eq:models-brier}
\end{equation}
按这里的求和定义，它位于$0$到$2$之间。相比交叉熵，它对极端错误的惩罚有上限，
因此两者会对少数非常自信的错误赋予不同权重。
Brier评分可以用来训练或评价概率预测，是较早用于概率预报的准则之一\scite{brier1950}

学习整个数据分布时，真实分布通常未知，只能得到样本。
设模型密度为$p_{\bm{\theta}}(\bm{x})$，常用平均负对数似然作为目标：
\begin{equation}
  \ell_{\mathrm{NLL}}(\bm{\theta})
  =-\frac{1}{n}\sum_{i=1}^{n}\log p_{\bm{\theta}}(\bm{x}_i).
  \label{eq:models-negative-log-likelihood}
\end{equation}
它要求模型对已观测数据给出较大的概率或密度，适用于GMM等分布模型。
在独立同分布采样、相关期望有限等条件下，样本平均估计总体的期望负对数似然；
最小化该总体目标等价于减小真实分布到模型分布的KL散度。
若建模的是给定输入后的类别分布$p_{\bm{\theta}}(y\mid\bm{x})$，其负对数似然
则对应上面的分类交叉熵。

交叉熵与KL散度在目标固定时给出相同的优化方向，Brier评分则改变了对极端错误的权重。
这些度量对不同类别的处理通常只依据概率大小，不自动考虑类别在空间上相距多远。
如果任务关心“概率质量移动到了多远的位置”，还需要将结果之间的距离纳入比较。

\subsection{关系与结构准则}
\label{sec:models-spatial-objectives}

逐样本差异要求输出接近给定目标，但有些任务没有这样的答案，或者还要约束多个样本、
参数与表示之间的关系。聚类希望同组样本聚在一起，检索希望相关样本比不相关样本更近；
分类可以在预测正确之外要求样本与边界留出间隔，降维则可以要求表示落在低维空间中。
这里把这些要求归为关系与结构准则，其中既有必须满足的限制，也有由损失鼓励的几何性质。

\paragraph{样本到中心的距离}

K-means用$k$个中心近似数据，最小化样本到最近中心的平方距离：
\begin{equation}
  \min_{\bm{\mu}_1,\ldots,\bm{\mu}_k}
  \frac{1}{n}\sum_{i=1}^{n}\min_{1\leq j\leq k}
  \lVert\bm{x}_i-\bm{\mu}_j\rVert_2^2.
  \label{eq:models-kmeans-objective}
\end{equation}
内层为每个样本选择中心，外层调整中心位置。这里没有真实类别作为参照，簇编号是
减小几何误差的结果。平方距离使较远的样本影响较大，也偏好用围绕中心的紧凑区域表示群组。
若改用$L_1$距离，给定分组时的中心变为各坐标的中位数，通常更能抵抗少量远离群组的点；
若主要关心方向，则可以使用基于余弦相似度的聚类目标。

标准K-means直接优化的是簇内紧凑程度，没有要求任意两个簇之间的距离都尽量大。
当中心取各簇均值时，固定数据的总平方离差可分成簇内与簇间两项，减小簇内项也会增大
簇间项；但后者描述各中心相对总体均值的分散程度，并不保证每对簇都具有大的分离间隔。

\paragraph{样本到分类边界的间隔}

SVM关心正确分类之外还能留出多少余量。对于$y\in\{-1,+1\}$和得分
$f(\bm{x})=\bm{w}^{\top}\bm{x}+b$，常用\term{合页损失}（hinge loss）为
\begin{equation}
  \ell_{\mathrm{hinge}}(y,f(\bm{x}))
  =\max\{0,1-yf(\bm{x})\}.
  \label{eq:models-hinge}
\end{equation}
当样本在正确一侧且得分达到规定余量时，损失为零；分类错误或余量不足时，损失增加。
但只把得分放大，也可能降低合页损失，几何边界却完全没动。线性SVM因此把这一损失与
$\lVert\bm{w}\rVert_2^2$的惩罚结合。对于$\bm{w}\ne\bm{0}$，样本的有符号几何
距离是$yf(\bm{x})/\lVert\bm{w}\rVert_2$；当得分余量固定为$1$时，较小的权重范数
对应较大的几何间隔。两项共同权衡间隔违反与间隔大小。
它利用类别标签组织空间，无需让同类样本都围绕同一个中心。

\paragraph{相关样本之间的相对距离}

检索与表示学习常只要求相对次序：同一物体的两张照片，应比不同物体的照片更接近。
设$\bm{a}$为待比较样本的表示，$\bm{p}$为相关样本的表示，$\bm{n}$为不相关样本的表示，
$d$为选定距离，$m>0$为所需余量。将一个参照样本、一个相关样本和一个不相关样本
放在一起比较，可以得到\term{三元组损失}（triplet loss）：
\begin{equation}
  \ell_{\mathrm{triplet}}
  =\max\{0,d(\bm{a},\bm{p})-d(\bm{a},\bm{n})+m\}.
  \label{eq:models-triplet}
\end{equation}
当相关样本至少比不相关样本近$m$时，损失为零；否则就按不足的差额计入损失。
它不规定所有表示必须落在某个固定中心附近。
这种目标适合学习距离关系，但效果依赖所选距离、样本配对和余量：过于容易的比较几乎
不提供学习信号，错误配对则可能把不该靠近的样本拉到一起。
图\ref{fig:models-spatial-objectives}对照了中心距离、分类间隔与相对距离三种要求。
利用已知相似关系学习表示或距离，是\term{度量学习}（metric learning）的基本思路\scite{weinberger2009}

\begin{figure*}[t]
  \centering
  % 原创教学几何示意；三幅图均使用固定坐标。
\begin{tikzpicture}[x=1mm,y=1mm,font=\footnotesize]
  \foreach \shift/\title in {0/靠近中心,57/留出分类间隔,114/比较相对远近}{
    \begin{scope}[xshift=\shift mm]
      \node[font=\heiti\small] at (22,44) {\title};
      \fill[BookPaper] (0,0) rectangle (46,36);
    \end{scope}
  }
  \foreach \x/\y in {5/6,8/17,16/7,20/14}{
    \draw[FigureInput!45] (\x,\y)--(12,11);
    \fill[FigureInput] (\x,\y) circle (1);
  }
  \foreach \x/\y in {26/22,31/31,40/25,38/18}{
    \draw[FigureLoss!50] (\x,\y)--(34,24);
    \fill[FigureLoss] (\x-0.9,\y-0.9) rectangle (\x+0.9,\y+0.9);
  }
  \foreach \x/\y in {12/11,34/24}{
    \draw[FigureModel,very thick] (\x-2,\y)--(\x+2,\y) (\x,\y-2)--(\x,\y+2);
  }
  \begin{scope}[xshift=57mm]
    \fill[FigureModel!10] (18,0) rectangle (28,36);
    \draw[FigureModel,thick] (23,0)--(23,36);
    \draw[FigureModel,dashed] (18,0)--(18,36) (28,0)--(28,36);
    \foreach \x/\y in {6/8,10/22,15/30,18/15}{
      \fill[FigureInput] (\x,\y) circle (1);
    }
    \foreach \x/\y in {28/9,32/25,37/17,40/31}{
      \fill[FigureLoss] (\x-0.9,\y-0.9) rectangle (\x+0.9,\y+0.9);
    }
    \draw[<->,BookMuted] (18,39)--(28,39);
  \end{scope}
  \begin{scope}[xshift=114mm]
    \draw[FigureModel,thick] (9,10)--(18,24);
    \draw[FigureLoss,thick,dashed] (9,10)--(37,17);
    \fill[FigureInput] (9,10) circle (1.4);
    \fill[FigureModel] (18,24) circle (1.4);
    \fill[FigureLoss] (36,16) rectangle (38,18);
    \node[below] at (9,8) {$\bm a$};
    \node[above] at (18,26) {$\bm p$};
    \node[above] at (37,19) {$\bm n$};
  \end{scope}
  \node at (23,-7) {样本与原型};
  \node at (80,-7) {样本与边界};
  \node at (137,-7) {样本与样本};
\end{tikzpicture}

  \caption{三种几何要求。左侧缩短样本到所属中心的距离；中间让带标签样本与分类边界保持
  间隔；右侧要求相关样本比不相关样本更接近参照样本。三者分别比较不同的几何对象。}
  \label{fig:models-spatial-objectives}
\end{figure*}

PCA则限制表示所在空间的维数，并在这个限制下减小重构误差。
这些例子说明，选择空间目标需要先明确希望保留什么：中心距离强调紧凑性，分类间隔
强调两类的分离，相对距离强调邻近次序，低维约束强调表示的压缩。
同一个模型也可以同时使用其中几种要求。

\FloatBarrier
\subsection{轨迹与长期目标}
\label{sec:models-sequence-objectives}

单次输出的误差还不能评价所有任务。机器人导航时，眼前距离终点更近的一步可能把它
带入死路，暂时绕行反而能更早到达。游戏决策和连续资源调度也有类似的问题：当前动作
会影响后续状态与奖励。评价这类模型，需要把整段行动过程计入，并考虑环境和动作选择
带来的多种可能经历。

一套根据当前信息选择动作的规则会与环境共同产生整段轨迹。常见准则先把一条轨迹中的
多步奖励合成总回报，再对各种可能轨迹取平均；也可以评价长期运行中的单位时间平均回报。
这些选择改变了哪些后果受到重视，不能由某一步的即时误差代替。第~\ref{sec:reinforcement-learning}
节将在建立状态、动作、奖励和策略等概念后，正式写出相应目标。

无论选择哪种回报，仅比较期望都可能掩盖结果的波动；若任务还要求控制失败概率，需要另加
相应约束。轨迹长、罕见事件多时，有限次交互对长期表现的估计也可能不稳定。

强化学习系统内部仍可以训练数值预测模型，例如让价值模型预测未来累计回报，
并用平方误差评价预测。这个局部损失服务于整体决策目标。
还要区分“输出具有序列形式”与“目标评价一整段决策过程”。例如，预测句子的下一个词
虽然产生文本序列，却可以逐位置比较类别概率，仍属于概率分布拟合。

从逐样本误差到整段轨迹的期望，上述目标都在说明“什么结果更好”；梯度下降、交替更新
等算法则说明“怎样改善结果”。
这种改善也可能只发生在局部：决策树逐步选择分裂条件，每一步都评价候选分裂，却未必
在优化一个关于整棵树的单一全局目标。近邻分类甚至可以直接保存样本，把主要计算留到
查询时。比较模型时，除了写出准则，还需要说明它用于整个模型，还是某一步局部选择。

\FloatBarrier
\section{模型维度的交叉与组合}
\label{sec:models-crossing}

前面的例子已经多次遇到维度交叉：一个非线性模型仍可以只有固定数量的参数，一个概率
模型也可以只返回一个数。现在把这些选择放回完整任务中，看看改变某一项设定，会保留
哪些性质，又会牵动哪些选择。

\subsection{模型维度的组合}
\label{sec:models-dimension-combination}

仍以根据气象特征预测明日温度为例。固定输入表示以后，一个完整的基准方案可以是：
用固定数量的权重计算输入的仿射函数，直接返回一个温度数值，并用平方误差训练。
沿五个维度描述，它是参数化、
非概率、线性结构的模型，输出接口返回标量，损失与学习准则衡量预测值与真实温度的差异。这些性质都
来自具体设定，任务名称本身只说明了我们要预测什么。

保留这个任务，改变预测的描述方式。若仍用仿射函数计算高斯分布的均值，
并额外估计一个方差参数，模型仍属固定维参数化模型，但表示规模增加；关系形式保持不变，
概率语义改为描述温度分布，损失与学习准则也相应改为评价观测数据的似然。

也可以保留点预测，改用含非线性激活的固定结构神经网络。此时参数化方式仍为参数化，
关系形式变为非线性；如果让网络输出分布参数，又得到非线性的概率预测。近邻回归
依赖不断增加的训练实例，可以形成非参数化的非概率模型；高斯过程同样不受固定有限维
系数限制，却对可能的函数和预测结果建立概率分布。

这些组合要求输出接口、关系形式和损失与学习准则相互配合。标量温度可以用绝对误差或平方误差
评价；类别分布可以用交叉熵或Brier评分评价；类别编号只表示身份，直接相减没有误差
意义。如果模型还预测方差，仅用均值的误差就无法评价它提供的全部信息。

关系形式也会影响损失与学习准则的选择。K-means用有限个中心表示群组，以样本到中心的距离衡量簇内
紧凑性；PCA把表示限制在低维子空间中，在这个约束下评价重构误差。把五个维度放在一起，
是为了检查参数化方式与表示规模、概率语义、关系形式、输出接口和损失与学习准则能否共同完成任务。

\subsection{分类模型的五维比较}
\label{sec:models-classification-grid}

把这些问题用于分类，可以检验章首提到的权重、规则树与近邻究竟差在哪里。
表\ref{tab:models-classifiers}列出九种具体设定。阅读时可以沿着一列比较同一属性，
也可以沿着一行检查一个完整模型的各项选择。

表中有两处设定需要说明。前文的高斯判别分析用各类的高斯分布计算类别后验；其中
\term{协方差}（covariance）描述各方向上的分散程度和相关性。各类共享同一个非奇异
协方差矩阵时，对数后验之差是线性的；各类使用不同协方差时，一般会得到二次边界，
称为\term{二次判别分析}（quadratic discriminant analysis）。朴素贝叶斯则采用特征
只取$0$或$1$的伯努利版本，例如某个词是否出现。在各特征条件概率严格位于零与一之间
时，其类别对数几率可以写成这些二值特征的线性函数。

\begin{table*}[htbp]
  \centering
  \footnotesize
  \setlength{\tabcolsep}{4pt}
  \renewcommand{\arraystretch}{1.25}
  \begin{tabularx}{\linewidth}{@{}>{\raggedright\arraybackslash}p{28mm}>{\raggedright\arraybackslash}p{16mm}>{\raggedright\arraybackslash}p{11mm}>{\raggedright\arraybackslash}p{23mm}>{\raggedright\arraybackslash}p{22mm}>{\raggedright\arraybackslash}X@{}}
    \toprule
    \textbf{具体模型设定} & \textbf{参数化方式与表示规模} & \textbf{概率语义} & \textbf{关系形式} & \textbf{输出接口} & \textbf{损失与学习准则} \\
    \midrule
    线性SVM & 参数化 & 非概率 & 线性得分 & 得分、类别 & 合页损失与权重惩罚 \\
    逻辑回归 & 参数化 & 概率 & 线性对数几率 & 类别概率 & 条件负对数似然，即分类交叉熵 \\
    高斯判别分析，共享协方差 & 参数化 & 概率 & 线性对数后验之差 & 类别概率 & 带标签数据的联合似然 \\
    伯努利朴素贝叶斯 & 参数化 & 概率 & 二值特征的线性对数几率 & 类别概率 & 类别比例与各类内特征的似然 \\
    高斯核SVM & 非参数化 & 非概率 & 非线性得分 & 得分、类别 & 特征空间中的间隔目标 \\
    二次判别分析 & 参数化 & 概率 & 一般为二次边界 & 类别概率 & 带标签数据的联合似然 \\
    固定结构的非线性神经分类器 & 参数化 & 概率 & 非线性得分 & 类别概率 & 通常使用分类交叉熵 \\
    硬投票$k$近邻 & 非参数化 & 非概率 & 局部投票形成的分段规则 & 类别 & 保存样本；查询时按距离选择邻居 \\
    硬标签决策树，叶数可增长 & 非参数化 & 非概率 & 条件分支与分段规则 & 类别 & 逐步选择使节点内类别更一致的分裂 \\
    \bottomrule
  \end{tabularx}
  \caption{分类模型的五维比较。表中采用所列的具体版本；概率模型均可进一步返回最大概率类别。
  “概率语义”按是否建立类别或联合分布判断，不能只由最终是否返回标签判断。}
  \label{tab:models-classifiers}
\end{table*}
\FloatBarrier

从“关系形式”一列看，线性SVM、逻辑回归、共享协方差的高斯判别分析和表中的
伯努利朴素贝叶斯都能得到线性得分或线性对数几率，而且都由预先确定数量的参数描述。
它们在关系形式、参数化方式与表示规模上相近，概率语义和损失与学习准则却不同。线性SVM不建立类别分布，而是用
间隔目标学习分类得分；逻辑回归直接建立$p(y\mid\bm{x})$，用条件负对数似然学习；
高斯判别分析和朴素贝叶斯则建立$p(y)p(\bm{x}\mid y)$，再由联合分布计算类别后验。
画出一条直线还不足以区分这些模型，还要看这条线来自怎样的概率假设与训练准则。

若线性边界不足以区分类别，以上几条路线又有不同的扩展方式。把线性SVM的线性核换成
高斯核，间隔准则和得分输出仍然保留，但边界可以弯曲，表示规模也可能随支持向量数量
增长。高斯判别分析在
各类共享协方差时产生线性边界，改为各类使用不同协方差后成为二次判别分析；它仍是参数化
概率模型，仍通过联合似然学习，只是边界通常变为二次曲面。逻辑回归与固定结构的非线性
神经分类器都可以直接输出$p(y\mid\bm{x})$并使用分类交叉熵，但后者通过中间层形成非线性
得分。这些扩展增加了非线性表达能力，却不一定改变参数化或非参数化的归类，概率语义也
可以保持原样。即使仍属参数化模型，候选关系的范围和容量也可能已经改变。

还可以从样本附近或条件区域内决定类别。硬投票$k$近邻保存训练样本，在查询
时按距离选择邻居；硬标签决策树逐步选择分裂条件，在叶节点给出类别。表中的两个版本都不
显式建立类别概率，表示规模也都可以随数据增长，但它们形成非线性的方式并不相同：近邻
方法产生局部的分段规则，决策树则把输入空间划分为若干条件区域。它们都能返回类别，学习
过程却分别侧重保存与查询规则、以及逐步改善节点内类别一致性。

现在沿输出接口一列回看：SVM的两个版本都能返回得分，多种关系形式也都能返回
类别概率，相同输出接口背后的参数化方式与表示规模、损失与学习准则却可能不同。表中的硬投票
与硬标签也是有意选定的版本；若进一步把近邻频率或叶节点类别比例解释为概率估计，概率语义
与输出接口两列就需要随之改写。方法名称
给出的是一组可能选择，表中的具体设定才决定本次比较的结论。

\FloatBarrier
\subsection{聚类模型的五维比较}
\label{sec:models-clustering-grid}

如果邮件没有标签，问题就从“哪一类是正确答案”变成“哪些邮件应当归在一起”。相似
可以指离同一个中心近，也可以指属于同一种分布，或由相邻样本连接成一片区域。
表\ref{tab:models-clusterers}列出七种具体设定；它们的差别始于这些分组依据，并进一步
影响表示规模与输出。

\begin{table*}[htbp]
  \centering
  \footnotesize
  \setlength{\tabcolsep}{4pt}
  \renewcommand{\arraystretch}{1.25}
  \begin{tabularx}{\linewidth}{@{}>{\raggedright\arraybackslash}p{28mm}>{\raggedright\arraybackslash}p{16mm}>{\raggedright\arraybackslash}p{11mm}>{\raggedright\arraybackslash}p{23mm}>{\raggedright\arraybackslash}p{22mm}>{\raggedright\arraybackslash}X@{}}
    \toprule
    \textbf{具体模型设定} & \textbf{参数化方式与表示规模} & \textbf{概率语义} & \textbf{关系形式} & \textbf{输出接口} & \textbf{损失与学习准则} \\
    \midrule
    K-means & 参数化 & 非概率 & 线性两两分配边界 & 簇编号 & 样本到中心的平方距离 \\
    高斯核K-means & 非参数化 & 非概率 & 非线性分配边界 & 簇编号 & 核特征空间中的平方紧凑性 \\
    模糊C均值 & 参数化 & 非概率 & 隶属度非线性；硬分界线性 & 隶属度向量、编号 & 隶属度加权的簇内平方距离 \\
    GMM，共享协方差 & 参数化 & 概率 & 线性两两分配边界 & 组别概率、编号 & 观测数据的似然 \\
    GMM，不同协方差 & 参数化 & 概率 & 一般为二次分配边界 & 组别概率、编号 & 观测数据的似然 \\
    DBSCAN & 非参数化 & 非概率 & 邻域连通关系 & 簇编号、噪声标记 & 邻域密度与连通规则，无统一的误差最小化目标 \\
    Ward层次聚类 & 非参数化 & 非概率 & 逐步合并的层次关系 & 层次树、截取后的簇编号 & 每步最小化簇内平方离差的增量 \\
    \bottomrule
  \end{tabularx}
  \caption{聚类方法的五维比较。参数化设定固定输入维数与中心或成分数。
  DBSCAN与层次聚类主要组织已有样本，其邻域和层次关系不直接对应一个线性预测函数。}
  \label{tab:models-clusterers}
\end{table*}
\FloatBarrier

标准K-means从中心出发定义群组，用有限个中心表示数据，以样本到中心的平方距离
衡量簇内紧凑性，并返回最近中心的编号；其两两分配边界是线性的，推导见
式\eqref{eq:models-kmeans-boundary}。高斯核K-means保留中心与紧凑性这一思路，但改在
核特征空间中比较样本与原型，因此原始空间中的边界可以弯曲。使用高斯核时，原型通常通过
训练样本表示，表示规模可以随数据增长；它与标准K-means都属于非概率模型，却分别采用
非参数化的非线性关系形式和参数化的线性分配关系\scite{dhillon2004}

处于群组交界处的样本，也可以同时与多个中心保持联系。\term{模糊C均值}
（fuzzy c-means）使用有限个中心和距离目标，为样本返回各簇的隶属度，也就是它在各簇中
参与计算的程度。隶属度非负且总和为一，其含义来自加权的簇内平方距离，
并不是由观测概率分布推断得到的组别后验。若只取最大隶属度对应的簇，分配又回到最近中心
判断。它与K-means在参数化方式与表示规模、概率语义上相近，主要区别落在输出接口和
损失与学习准则上
\scite{bezdek1984}

如果希望归属程度具有概率含义，就需要为数据建立相应的分布。GMM以多个高斯分布及其
混合比例建立观测数据与潜在组别的联合分布，通过后验概率表示样本属于各组的可能性，
并用数据似然学习参数。各成分共享同一协方差
时，两成分对数后验之差中的二次项抵消，得到线性分配边界；各成分使用不同协方差时，二次项
通常保留，边界一般为二次曲面。两个版本都是参数化概率模型，输出接口也都返回组别概率，
关系形式的选择
却不同。共享协方差GMM与K-means都能形成线性两两边界，但前者给出概率解释并优化似然，
后者只按几何距离分组。这一对比说明，线性关系形式本身不能决定模型是否具有概率语义。

当群组沿弯曲路径展开时，单个中心或高斯成分可能难以描述形状。\term{DBSCAN}
（density-based spatial clustering of applications with noise）根据局部邻域中的样本数量寻找稠密
区域，再按连通关系扩展群组；它可以形成弯曲群组，并把孤立样本标为噪声。这里的“密度”
是邻域计数，不是归一化的观测概率密度，因此表中的DBSCAN是非参数化、非概率模型，输出
包括簇编号与噪声标记\scite{ester1996}
如果希望同时保留不同粗细程度的分组，可以使用\term{凝聚层次聚类}
（agglomerative hierarchical clustering）：从小群组开始逐步合并，
保留一棵可以在不同层级截取的树。表中的\term{Ward方法}每一步选择使簇内平方离差增加
最少的合并；它与K-means都重视平方紧凑性，却输出层次关系而不是一组固定中心
\scite{ward1963}

这些比较给出了几种具体取舍：希望弯曲的分配边界，可以从K-means改用核表示；希望保留
多重归属，可以考虑隶属度或组别概率，并分清二者的含义；希望保留局部连通形状或不同
层级，则需要邻域关系或合并树。即使最终都截取为一组簇编号，各方法保留的信息和认定
“应当归在一起”的理由仍然不同。

\FloatBarrier
\subsection{模型组合与选择}
\label{sec:models-combination-selection}

前两个小节分析的是单个模型的维度组合，实际系统还可能把多个模型或表示模块连接起来。
这时除了判断各模块的属性，还要考察它们怎样协作。
例如，先用PCA压缩邮件特征，再用K-means分组，是把表示与聚类串联；
先用神经网络学习表示，再接一个线性分类器，是把非线性特征计算与线性得分结合起来。
后一个模型虽然最后一层是线性的，整体对原始输入通常仍然是非线性的。

连接模块还会带来训练上的选择。分阶段学习先训练上游表示，再固定它并训练下游模型；
端到端学习根据最终目标共同调整相关参数。例如，压缩模块即使重构误差很小，也可能
丢掉对分类有用的微小差别。模块自己的目标是否保留了后续任务需要的信息，需要单独检查。

多个模型也可以共同完成同一项预测，这种做法称为\term{集成学习}（ensemble learning）。
例如，随机森林汇集多棵树的结果，梯度提升逐步加入模型以改善既有预测。此时需要说明
各个基模型怎样产生差异，以及它们的输出按什么规则合并。

无论选择单个模型还是组合模型，五维分析都只是明确假设，并不能代替实际比较。最终选择
还需要结合数据、任务与资源约束，在相同验证过程下比较预测表现，同时考察存储和查询成本、
解释需求，以及面对分布变化时的可靠性。

\section{本章小结}
\label{sec:models-summary}

回到章首的邮件问题：学习权重、建立规则树和保存近邻，都能从标签中学习，却把规律保存
成了不同形式。模型族规定了允许作哪些选择，参数或实例表示确定了实际使用的关系，
学习算法则负责根据数据完成选择。

五个维度让这些差别可以逐项比较。参数化方式与表示规模说明表示规模受到什么限制，概率语义
说明是否用分布描述变量，关系形式说明如何连接变量，输出接口说明结果怎样组织以及代表什么，
损失与学习准则说明评价针对逐样本目标、多个对象之间的关系，还是整段轨迹的长期结果。
它们需要分开判断，也需要相互配合：线性与非线性结构都能用于概率预测，固定数量的参数
也能确定一个分布；类别编号与簇编号虽然形式相同，却需要不同的学习依据。

\paragraph{延伸阅读}

模型说明了怎样表示规律，第\ref{chap:machine-learning-paradigms}章将进一步讨论经验和反馈
改变时，怎样组织学习。模型在新数据上是否可靠，还要靠验证与进一步的理论分析。
第\ref{part:learning-theory}部分将讨论有限经验在什么条件下能够支持模型泛化；此后，第\ref{part:classic-models}、
第\ref{part:neural-network-models}和第\ref{part:probabilistic-models}部分将分别展开经典模型、
神经网络模型与概率图模型。
